%! Unit 1 cover sheet — the shared body.
%
% Authored ONCE and \input by both wrappers, so the student cover and the key
% cover carry a byte-identical page 1 and cannot drift:
%   unit_cover/main.tex      → this body alone (1pp, merged into the student packet)
%   unit_cover_key/main.tex  → this body + the exam-scoring teacher notes on p2
%                              (2pp, merged into the key packet only)
%
% unit.mk merges this sheet into the unit packet WITHOUT the packet-wide
% pagination pass (that runs per lesson), so a folio here would print a stray
% "1" ahead of Lesson 1.0's own page 1. The unit cover carries no number.
\pagestyle{empty}

% ── Full-bleed burgundy banner ───────────────────────────────────────────────
% The header text is set INSIDE a tikz node anchored to the page so it can never
% drift above the paper edge (a negative \vspace here pushes it off the sheet).
\begin{tikzpicture}[remember picture, overlay]
  \fill[burgundy] (current page.north west)
    rectangle ([yshift=-1.70in]current page.north east);
  \node[anchor=north, align=center, inner sep=0pt, text width=\paperwidth]
    at ([yshift=-0.32in]current page.north) {%
      {\color{white}\bfseries\fontsize{30}{34}\selectfont Linear Algebra}\\[11pt]
      {\color{blushmid}\bfseries\fontsize{19}{23}\selectfont Unit 1: Vectors and Matrices}\\[11pt]%
    };
\end{tikzpicture}

% Push the body clear of the 1.70in banner (page top margin is 0.60in).
\vspace*{1.32in}

% ── Unit overview ────────────────────────────────────────────────────────────
\begin{tcolorbox}[colback=blush,colframe=burgundy,boxrule=0.9pt,arc=2mm,
    left=4mm,right=4mm,top=3mm,bottom=3mm]
  \textbf{Unit Overview.}\quad
  Unit 1 builds the whole course out of two operations: \emph{scaling} a vector and
  \emph{adding} vectors. Students meet a vector as both a list of numbers and an arrow in
  the plane, then combine vectors ($c\,\vv + d\,\ww$) and ask the question that runs through
  the rest of the course --- \emph{which points can these vectors reach?} The dot product
  turns that algebra into geometry, giving length, angle, and a test for perpendicularity
  from a single number. The unit closes by standing vectors side by side into a
  \textbf{matrix}, reading $A\xx$ as a combination of the columns of $A$, and factoring
  $A = CR$ to separate a matrix's independent columns from the recipe that rebuilds the rest.
\end{tcolorbox}

\vspace{0.13in}

% ── Lessons in this unit ─────────────────────────────────────────────────────
\renewcommand{\arraystretch}{1.32}
\begin{skillbox}[Lessons in This Unit]{goldbox}
\begin{tabularx}{\linewidth}{@{} c >{\bfseries\raggedright\arraybackslash}p{0.29\linewidth} X @{}}
  \textbf{\#} & \textbf{Title} & \textbf{Focus} \\
  \hline
  1.0 & Introducing Vectors: Vocabulary and Length
    & Components and arrows; adding and scaling and what each does to the picture;
      length $\|\vv\| = \sqrt{v_1^{\,2}+v_2^{\,2}}$ as straight-line distance \\
  \hline
  1.1 & Linear Combinations of Vectors
    & $c\,\vv + d\,\ww$ as scale-then-add; the multiples of one vector fill a line and
      two non-parallel vectors fill the plane; finding the weights for a target \\
  \hline
  1.2 & Lengths and Angles from Dot Products
    & $\vv\cdot\ww$ and its geometry: $\|\vv\| = \sqrt{\vv\cdot\vv}$, the angle from
      $\cos\theta = \frac{\vv\cdot\ww}{\|\vv\|\,\|\ww\|}$, and $\vv\cdot\ww = 0$ as
      perpendicular \\
  \hline
  1.3 & Matrices and Column Spaces
    & A matrix as columns standing side by side; $A\xx$ as a combination of those columns;
      the column space $C(A)$ --- every target $\bb$ the columns can build \\
  \hline
  1.4 & Matrix Multiplication and $A = CR$
    & Multiplying matrices \emph{by columns}; factoring $A = CR$ with $C$ the independent
      columns and $R$ the recipe; \textbf{rank} as the count of independent columns \\
  \hline
\end{tabularx}
\end{skillbox}

\vspace{0.11in}

% ── Big ideas of the unit ────────────────────────────────────────────────────
\begin{skillbox}[Big Ideas of Unit 1]{greenbox}
\small
\begin{itemize}[leftmargin=1.3em, itemsep=3pt, topsep=1pt]
  \item A vector is a list of numbers \emph{and} an arrow. The algebra and the geometry are
        two views of the same object --- every computation should have a picture.
  \item \textbf{Scaling} and \textbf{adding} generate everything here. A linear combination
        does both at once, and the set of points it can reach --- the \textbf{span} --- has a
        shape: a line, a plane, or all of space.
  \item The \textbf{dot product} is one number that carries geometry: length, angle, and the
        perpendicularity test $\vv\cdot\ww = 0$.
  \item $A\xx$ is a \textbf{linear combination of the columns of $A$}. So ``does $A\xx = \bb$
        have a solution?'' is the same question as ``is $\bb$ in the column space $C(A)$?''
  \item \textbf{Rank} counts the independent columns. $A = CR$ makes that count visible: $C$
        keeps the independent columns, $R$ says how to rebuild the others.
\end{itemize}
\end{skillbox}

\vspace{0.10in}

\begin{remindbox}
\small
\textbf{Source.} Unit 1 follows Chapter 1 of Gilbert Strang's \emph{Linear Algebra for
Everyone}, sections 1.1--1.4.\quad
\textbf{Assessment.} The unit ends with a test; a full \textbf{practice test} is bound at the
back of this packet.
\end{remindbox}
